\answer{ex:12:1}{$0.60$ (без кореляції було б $0.18$), межа $\sqrt{c/(rT)}\approx0.58$: сто нейронів розрізняють лише «є/немає».}
\answer{ex:12:2}{(а)~$10^8$ імпульсів, $\approx10\,\text{мДж}$: $0.3\,\%$ від $3$~Дж усього мозку. (б)~$3$~Дж, у $\approx300$ разів більше.}
\answer{ex:12:3}{(а)~$\theta_A\in[2,3)$, $\theta_B\in[1,2)$; чужа фігура дає не більше за $2$ і $1$. (б)~$10\cdot96^2\approx9.2\cdot10^4$.}
\answer{ex:12:4}{(а)~$T_d=0.85\,\tau_a=0.34\,\text{с}$. (б)~$T_d\propto\tau_a$ і не залежить від $I$: $\tau_a\approx3.5\,\text{с}$ (моделювання: $3.1\,\text{с}$).}
\answer{ex:12:5}{$p\approx0.17$. Частка падає повільно: $0.90$, $0.88$, $0.86$, $0.84$, $0.82$, $0.79$ при $N=8$, $12$, $24$, $48$, $96$, $192$.}
\answer{ex:12:6}{Усі десять прогонів безпомилкові при $\tau_{\text{сл}}=20$, $50$ і $200\ms$ (при $30\ms$ один прогін із десяти зірвався); при $5$--$10\ms$ --- близько $0.5$, при $0.5$--$2\,\text{с}$ якість падає до $0.86$. Згори: слід має згаснути за паузу, $e^{-0.3\,\text{с}/\tau_{\text{сл}}}\ll1$.}
\answer{ex:12:7}{Однієї затримки замало: вікно $\pm5\ms$ не покриває $\Delta x/v$ від $5$ до $20\ms$. Дві гілки гальмування з $5<d_1\le10\ms$ і $15\le d_2\le d_1+10\ms$.}
\answer{ex:12:8}{\texttt{two\_stage}: $\Delta=2$, нічия від одного перевертання, зміна відповіді --- від двох. Лічильник одиниць: одне перевертання; $0.57$ при $p=0.1$.}
